\documentclass[11pt,a4paper]{article}
\usepackage[T1]{fontenc}
\usepackage{lmodern}
\usepackage[margin=27mm]{geometry}
\usepackage{amsmath,amssymb,amsthm,booktabs,microtype}
\usepackage[hidelinks]{hyperref}
\newtheorem{theorem}{Theorem}
\newtheorem{lemma}[theorem]{Lemma}
\newtheorem{proposition}[theorem]{Proposition}
\newtheorem{corollary}[theorem]{Corollary}
\theoremstyle{definition}
\newtheorem{definition}[theorem]{Definition}
\theoremstyle{remark}
\newtheorem{remark}[theorem]{Remark}
\newcommand{\Ostar}{O^{*}}
\newcommand{\E}{\mathbb E}
\newcommand{\Prb}{\mathbb P}
\newcommand{\one}{\mathbf 1}
\emergencystretch=2em
\clubpenalty=10000
\widowpenalty=10000
\hypersetup{pdftitle={Weighted and Round-Based Early Enumeration}}
\title{Weighted and Round-Based Early Enumeration\\
\large Random signs on every graph, Max-Cut on spectral expanders,\\
and count guessing without a given decomposition}
\author{}
\date{12 October 2026}
\begin{document}
\maketitle

\begin{abstract}
Two modifications of early enumeration remove the main restrictions of the
earlier class theorems. Weighting each learned constraint by the inverse square
root of the degree gives a fixed exponential saving for Max 2-SAT with random
literal signs on \emph{every} interaction graph, with arbitrary unit clauses and
several clauses per edge allowed; the earlier bound required a bounded ratio of
maximum to minimum degree. Fixing the preferred values of the branching
variables only between a constant number of rounds makes the analysis depend on
the ordinary spectral norm of the remainder instead of the norm of its
triangular part. As a consequence, Max-Cut with arbitrary unit clauses is solved
in time $\Ostar(c^n)$, $c<2$, and polynomial space on every $d$-regular graph
whose nontrivial eigenvalues are $O(\sqrt d)$, for all $d_0\le d\le(1-\varepsilon)n$.
This includes Paley graphs, random regular graphs and Ramanujan graphs.
Finally, the counts of count guessing can be replaced by projections onto
eigenvectors, guessed on a polynomial grid; no decomposition has to be given.
All algorithms are deterministic, exact on every input, and use polynomial space.
\end{abstract}

\section{Results}\label{sec:results}

Write the score of a formula as
$F(x)=c+\sum_uh_ux_u+\sum_{u<v}J_{uv}x_ux_v$ on $x\in\{0,1\}^n$, and
$\mathrm{gain}_t(x)=F(x^{t\to1})-F(x^{t\to0})=h_t+\sum_{w\ne t}J_{tw}x_w$.
An equal-polarity clause, $(u\lor v)$ or $(\neg u\lor\neg v)$, contributes $-1$
to $J_{uv}$, a mixed one $+1$.

\begin{itemize}
\item \textbf{Random signs on every graph (Theorem~\ref{thm:weighted}).} Fix
any graph $G$ and any multiset of unit clauses. Put on each edge a random
multiset of at most $m_0$ clauses whose law is invariant under complementing
either endpoint, for example one clause with uniform literal signs, or each of
the four clauses independently with probability $q$. Weighted early enumeration
enumerates at most $3\cdot2^ne^{-c_0n}$ assignments in expectation, with $c_0>0$
depending only on $m_0$ and a non-degeneracy constant $\kappa$, not on $G$.
\item \textbf{Round-based enumeration (Theorem~\ref{thm:rounds}).} The class of
the count-guessing note, with its triangular-norm condition replaced by
$\|N\|^2\le K\,D_N$ for the ordinary operator norm of the remainder, and with
the aligned part an arbitrary low-rank matrix $\sum_k\lambda_kv_kv_k^{\mathsf T}$.
\item \textbf{Max-Cut on spectral expanders (Corollary~\ref{cor:maxcut}).}
Every $d$-regular graph with all nontrivial eigenvalues at most $\lambda_0\sqrt d$
in absolute value, $d_0(\lambda_0)\le d\le(1-\varepsilon)n$, with both clauses
$(u\lor v),(\neg u\lor\neg v)$ on each edge and arbitrary unit clauses: time
$\Ostar(c^n)$ with $c=c(\lambda_0,\varepsilon)<2$.
\item \textbf{Spectral count guessing (Corollary~\ref{cor:spectral},
Proposition~\ref{prop:planted}).} Guessing the projections of $x$ onto the top
eigenvectors of $J$ on a grid of polynomial size handles planted low-rank parity
structure without knowing it.
\end{itemize}

\section{Preliminaries}\label{sec:prelim}

\paragraph{Enumeration.} All algorithms branch on a prefix of $T$ variables,
learn affine constraints on the way, and at each prefix enumerate the
completions of one summed affine inequality. Such solutions can be listed with
polynomial delay, and the cost is $\Ostar(2^T+W)$, where $W$ is the number of
complete assignments satisfying the summed inequality of their own
prefix~\cite{EE,Paper}. For uniform $x$, $W=2^n\Prb[x\text{ passes}]$.

\paragraph{Probabilistic tools.}
\begin{itemize}
\item[(T1)] \emph{Azuma--Hoeffding.} If $X_1,X_2,\dots$ are martingale
differences and $X_i$ lies, given the past, in an interval of length $\ell_i$,
then $\Prb[\sum_iX_i\le-a]\le e^{-2a^2/\sum_i\ell_i^2}$.
\item[(T2)] \emph{Exponential supermartingale.} If, given the past, $X_i$ is a
Rademacher sum with variance proxy $V_i$ (so $\E e^{\lambda X_i}\le
e^{\lambda^2V_i/2}$) and $\sum_iV_i\le V$ always, then
$\Prb[\sum_iX_i\le-a]\le e^{-a^2/(2V)}$.
\item[(T3)] \emph{Talagrand.} If $f$ is convex and $\ell$-Lipschitz on
$[-1,1]^T$ and $u$ is uniform on $\{\pm1\}^T$, then
$\Prb[f\le\E f-r-8\sqrt\pi\,\ell]\le4e^{-r^2/(16\ell^2)}$~\cite{Tal95}.
\item[(T4)] \emph{Hanson--Wright.} There is an absolute constant $c_{\rm HW}>0$
such that for a symmetric matrix $M$ with zero diagonal and a uniform
$\beta\in\{\pm1\}^m$,
$\Prb[\beta^{\mathsf T}M\beta\ge s]\le\exp(-c_{\rm HW}\min\{s^2/\|M\|_F^2,\,
s/\|M\|\})$~\cite{RV13}.
\item[(T5)] \emph{Fourth moment.} A Rademacher sum $Z$ satisfies
$\E|Z|\ge\sqrt{\operatorname{Var}Z/3}$.
\end{itemize}

\begin{lemma}[Symmetric shift]\label{lem:shift}
If $Z$ is symmetric, $a\in\mathbb R$ and $K>0$, then
$\E\min\{|a+Z|,K\}\ge\frac12\E\min\{|Z|,K\}$.
\end{lemma}
\begin{proof}
The function $\phi(z)=\min\{z,K\}$ is concave and nondecreasing on $[0,\infty)$
with $\phi(0)=0$, hence subadditive. So $\phi(|a+Z|)+\phi(|a-Z|)\ge
\phi(|a+Z|+|a-Z|)\ge\phi(2|Z|)\ge\phi(|Z|)$. Take expectations and use that
$a+Z$ and $a-Z$ have the same law.
\end{proof}

\section{Random signs on every graph}\label{sec:weighted}

\begin{definition}[Complementation-invariant random signs]\label{def:invariant}
Let $G=(V,E)$ be a graph without isolated vertices, $n=|V|$, with any fixed order
$1,\dots,n$ of $V$, and let $T=\lceil n/2\rceil$. For each edge $e=\{u,v\}$ let
$M_e$ be a random multiset of binary clauses on $u,v$, independently over the
edges, such that
\begin{enumerate}
\item[(i)] $|M_e|\le m_0$;
\item[(ii)] the law of $M_e$ is invariant under replacing the literal of $u$ by
its complement in every clause of $M_e$, and likewise for $v$;
\item[(iii)] $\E\bigl[(\sum_{c\in M_e}s_u(c))^2\bigr]\ge\kappa>0$ for every edge
$e$ and endpoint $u$, where $s_u(c)\in\{\pm1\}$ is the sign of the literal of
$u$ in $c$.
\end{enumerate}
The unit clauses form an arbitrary multiset, fixed independently of the $M_e$.
\end{definition}

One clause per edge with independent uniform literal signs has $m_0=\kappa=1$.
Including each of the four clauses on every edge independently with probability
$q$ has $m_0=4$ and $\kappa=4q(1-q)$; with $q=\frac14$, this is the model in
which the four clauses of each pair of a fixed set are sampled with probability
$\frac14$.

\paragraph{Algorithm A (weighted early enumeration).}
For a clause $c\in M_{tv}$, the \emph{partner literal} is its literal on $v$.
Let $b_t$ be the number of unit clauses $(t)$ minus the number of unit clauses
$(\neg t)$, and $G_{tv}(x_v)=\sum\{s_t(c):c\in M_{tv},\ \text{the partner
literal of $c$ is false under }x_v\}$. Then
$\mathrm{gain}_t(x)=b_t+\sum_{v\sim t}G_{tv}(x_v)$. Branch on $x_1,\dots,x_T$.
At step $t$ let
\[
 \Delta_t=b_t+\sum_{v<t,\,v\sim t}G_{tv}(x_v)
  +\frac12\sum_{v>t,\,v\sim t}\ \sum_{c\in M_{tv}}s_t(c),
\]
the expected gain of $x_t$ when the later variables are uniform, and
$p_t=[\Delta_t\ge0]$. If $x_t\ne p_t$, then $t$ is a \emph{creator}. At each
prefix of length $T$ enumerate the completions satisfying
\begin{equation}\label{eq:weighted}
 \sum_{t\ \text{creator}}\omega_t\bigl((2p_t-1)\,\mathrm{gain}_t(x)+1\bigr)\le0,
 \qquad \omega_t=\deg_G(t)^{-1/2},
\end{equation}
and evaluate $F$ on each.

\begin{lemma}[Exactness]\label{lem:exactA}
Among the optima, let $x^\circ$ maximize, lexicographically, the vector of
indicators $[x_t=p_t]$, $t\le T$, along its own path. Then
$(2p_t-1)\mathrm{gain}_t(x^\circ)\le-1$ for every creator $t$ on its path, so
$x^\circ$ satisfies~\eqref{eq:weighted} and is enumerated.
\end{lemma}
\begin{proof}
Scores are integers. If $(2p_t-1)\mathrm{gain}_t(x^\circ)\ge0$ for a creator $t$,
flipping $x^\circ_t$ to $p_t$ gives an optimum with the same path before $t$ and
indicator one at $t$, which is lexicographically larger.
\end{proof}

\begin{theorem}[Random signs on every graph]\label{thm:weighted}
In the model of Definition~\ref{def:invariant}, Algorithm~A is exact on every
input, and for every graph, every order and every multiset of unit clauses the
expected number of enumerated assignments is at most $3\cdot2^ne^{-c_0n}$, where
\[
 c_0=\frac{c_1^2}{2704\,m_0^2},\qquad
 c_1=\frac{\sqrt\kappa}{2\sqrt{1+4m_0^4/\kappa^2}}.
\]
Hence the expected running time is $\Ostar(2^{n/2}+2^ne^{-c_0n})$, and the
running time exceeds $\Ostar(2^ne^{-c_0n/2})$ with probability at most
$3e^{-c_0n/2}$.
\end{theorem}

For one uniformly signed clause per edge, $c_0=1/54080$. The earlier bound for
this model, Theorem~4.1 of~\cite{Paper}, has base $2e^{-d/(1250D)}$ with $d,D$
the minimum and maximum degree, requires $d\ge6$, and excludes unit clauses.

\begin{proof}
\emph{Step 0: race identity.} Let $x$ be uniform and independent of the formula.
Write $\mathrm{gain}_t(x)=\Delta_t+\sum_{v>t,v\sim t}\xi_{tv}$ with
$\xi_{tv}=G_{tv}(x_v)-\frac12\sum_{c\in M_{tv}}s_t(c)$. For a creator,
$(2p_t-1)\Delta_t=|\Delta_t|$. Hence $x$ passes~\eqref{eq:weighted} if and only if
$S+Q\le0$, where, with $I$ the set of creators,
\[
 S=\sum_{t\in I}\omega_t(|\Delta_t|+1),\qquad
 Q=\sum_{t\in I}\omega_t(2p_t-1)\sum_{v>t,\,v\sim t}\xi_{tv}.
\]

\emph{Step 1: late signs.} By (ii), $M_e$ has the law of a multiset obtained from
an independent copy $\sigma_e$ of $M_e$, the \emph{skeleton}, by multiplying the
signs of $u$ in all its clauses by an independent fair sign $\rho_{u,e}$, and
those of $v$ by $\rho_{v,e}$. Put $\theta_{u,e}=\rho_{u,e}(2x_u-1)$. The literal of
$u$ in $c$ is false under $x_u$ exactly when $\sigma_u(c)\theta_{u,e}=-1$. All
$\theta_{u,e}$ are independent fair signs, independent of the skeletons and of
$x$, because the $\rho$'s are. Reveal the skeletons first; then, for
$v=1,\dots,n$, reveal the signs $\rho_{v,e}$ ($e\ni v$) and $x_v$. Let
$\mathcal F_v$ be the information after step $v$.

For $e=\{t,v\}$ we have $G_{tv}(x_v)=\rho_{t,e}\,a_e$ with
$a_e=\sum\{\sigma_t(c):\sigma_v(c)\theta_{v,e}=-1\}$. Hence
$\Delta_t=b_t+Z_t$, where
\[
 Z_t=\sum_{v<t}\rho_{t,e}\,a_e+\sum_{v>t}\rho_{t,e}\,a'_e,\qquad
 a'_e=\tfrac12\textstyle\sum_{c\in M_e}\sigma_t(c),
\]
and the coefficients $a_e,a'_e$ are $\mathcal F_{t-1}$-measurable. Given
$\mathcal F_{t-1}$, $Z_t$ is a Rademacher sum with variance
$Y_t=\sum_{v<t}a_e^2+\sum_{v>t}a_e'^2\le m_0^2\deg(t)$, and $x_t$ is uniform and
independent of $\Delta_t$, so $t$ is a creator with probability $\frac12$. For
$v>t$, $\xi_{tv}=\rho_{t,e}\,\theta_{v,e}\,\eta_e$ with
$\eta_e=\sum\{\sigma_t(c):\sigma_v(c)=-1\}-a'_e$ and $|\eta_e|\le m_0/2$.

\emph{Step 2: the shortfall.} Let
$\tilde Y_t=[t\in I]\,\omega_t\min\{|\Delta_t|,m_0\sqrt{\deg t}\}\in[0,m_0]$; then
$S\ge\sum_{t\le T}\tilde Y_t$. By Lemma~\ref{lem:shift}, (T5) and
$(z-K)_+\le z^2/(4K)$,
\[
 \E[\tilde Y_t\mid\mathcal F_{t-1}]
 \ge\tfrac12\omega_t\cdot\tfrac12\Bigl(\sqrt{Y_t/3}-\frac{Y_t}{4m_0\sqrt{\deg t}}\Bigr)
 \ge\tfrac14\omega_t\sqrt{Y_t}\Bigl(\tfrac1{\sqrt3}-\tfrac14\Bigr)
 \ge\mu_t:=\frac{\omega_t\sqrt{Y_t}}{13},
\]
using $Y_t\le m_0^2\deg t$. By (T1),
$\Prb[\sum_{t\le T}\tilde Y_t<\sum_{t\le T}\mu_t-a]\le e^{-2a^2/(Tm_0^2)}$.

\emph{Step 3: the mean shortfall.} For $v<t$, $a_e$ takes the values
$\alpha_\pm=\sum\{\sigma_t(c):\sigma_v(c)=\mp1\}$ with probability $\frac12$
each, so $\E a_e^2\ge\frac14\E(\alpha_++\alpha_-)^2
=\frac14\E(\sum_c\sigma_t(c))^2\ge\kappa/4$; also $\E a_e'^2\ge\kappa/4$. The
summands of $Y_t$ are independent over the edges and lie in $[0,m_0^2]$, so
$\E Y_t\ge\kappa\deg(t)/4$ and $\operatorname{Var}Y_t\le m_0^4\deg(t)/4$. Hence
\[
 \E\sqrt{Y_t}\ge\frac{(\E Y_t)^{3/2}}{(\E Y_t^2)^{1/2}}
 =\frac{\sqrt{\E Y_t}}{\sqrt{1+\operatorname{Var}Y_t/(\E Y_t)^2}}
 \ge c_1\sqrt{\deg t},
\]
and $\E\sum_{t\le T}\mu_t\ge c_1T/13$. The sum $\sum_{t\le T}\mu_t$ is a function
of the independent skeletons and signs $\theta$; those of one edge change at most
two terms, each by at most $\omega m_0/13$. Since $\sum_e(\omega_u+\omega_v)^2\le
2\sum_v\omega_v^2\deg v=2n$, McDiarmid's inequality gives
$\Prb[\sum_{t\le T}\mu_t<c_1T/13-b]\le e^{-169b^2/(m_0^2n)}$.

\emph{Step 4: the interaction term.} Order the summands of $Q$ by their later
endpoint $v$. The summand for the edge $\{t,v\}$ is
$[t\in I]\,\omega_t(2p_t-1)\rho_{t,e}\eta_e\cdot\theta_{v,e}$; its coefficient is
known before step $v$ and $\theta_{v,e}$ is a fresh fair sign, independent of the
other signs revealed at step $v$. By (T1), with interval lengths
$\omega_tm_0$ and $\sum_{t\le T}\omega_t^2\deg(t)\le T$,
$\Prb[Q\le-c]\le e^{-2c^2/(m_0^2T)}$.

\emph{Step 5.} Take $a=b=c=c_1T/52$. Outside the three exceptional events,
$S+Q\ge c_1T/13-3c_1T/52=c_1T/52>0$. With $T\ge n/2$ the three exponents are at
least $c_1^2T/(1352m_0^2)\ge c_0n$, so the three probabilities are at most
$e^{-c_0n}$ each, and
$\E W\le3\cdot2^ne^{-c_0n}$. Markov's inequality gives the tail statement.
\end{proof}

\begin{corollary}[Arbitrary supports]\label{cor:support}
Fix any set of variable pairs and include each of the four clauses on each of
these pairs independently with probability $q\in(0,1)$; add arbitrary unit
clauses (variables in no pair are set by their unit clauses). Then Max 2-SAT is solved in expected time $\Ostar(2^n e^{-c_0(q)n})$ with
$c_0(q)>0$ depending only on $q$. No condition on the density, the degree
sequence or the order is needed.
\end{corollary}

\begin{remark}[Complementary pairs]
Condition (iii) fails when every edge carries a complementary pair, as in Max-Cut
with random renamings (a $\pm J$ spin glass). Then the later-neighbour part of
$Y_t$ vanishes and the shortfall comes from earlier neighbours only. With a
uniformly random order, $\E\sqrt{Y_t}$ is of order $\sqrt{(t/n)\deg t}$ and the
same proof should go through; we do not write it out.
\end{remark}

\section{Round-based enumeration with guessed projections}\label{sec:rounds}

\paragraph{Setting.} Let $P=\sum_{k=1}^r\lambda_kv_kv_k^{\mathsf T}$ with
$\lambda_k\in\mathbb R$ and $v_k\in\mathbb R^n$, and let $N$ be the off-diagonal
part of $J-P$, so $N_{uu}=0$ and $J_{uv}=P_{uv}+N_{uv}$ for $u\ne v$. Put
$y_k(x)=\langle v_k,x\rangle$. For each $k$ fix a finite grid
$\mathcal Y_k\subset\mathbb R$ such that every $y_k(x)$, $x\in\{0,1\}^n$, is within
$\delta_k\ge0$ of a grid point, and let
\[
 s_t=\sum_k|\lambda_k|\,|v_k(t)|\,\bigl(\delta_k+|v_k(t)|\bigr).
\]
Let $T=\lfloor\tau n\rfloor$, and split the prefix into $\ell\ge2$ consecutive
\emph{rounds} $B_1,\dots,B_\ell$ of the branching order; the remaining
variables form round $\ell+1$. Write $\rho(t)$ for the round of $t$.

\paragraph{Algorithm C.} For each $\hat y\in\mathcal Y=\prod_k\mathcal Y_k$ run
the following. Put $h^{\hat y}=h+\sum_k\lambda_k\hat y_kv_k$ and
$F'_{\hat y}(x)=c+\langle h^{\hat y},x\rangle+\sum_{u<v}N_{uv}x_ux_v$. Branch on
the prefix. For $t\in B_j$ let
\[
 \Delta_t=h^{\hat y}_t+\sum_{s:\rho(s)<j}N_{ts}x_s
          +\frac12\sum_{w\ne t:\rho(w)\ge j}N_{tw}
\]
and $p_t=[\Delta_t\ge0]$; the preferred value of $t$ depends only on the earlier
rounds. If $x_t\ne p_t$, then $t$ is a creator and
$\gamma_t(x)=(2p_t-1)\bigl(F'_{\hat y}(x^{t\to1})-F'_{\hat y}(x^{t\to0})\bigr)$.
At each prefix enumerate the completions of
\begin{equation}\label{eq:roundsum}
 \sum_{t\ \text{creator}}\bigl(\gamma_t(x)-s_t\bigr)\le0,
\end{equation}
evaluate $F$ on them, and return the best assignment over all runs.

\begin{lemma}[Flip identity]\label{lem:flip}
If $|y_k(x)-\hat y_k|\le\delta_k$ for all $k$, then for every $t$ the gains of
$x_t$ under $F$ and under $F'_{\hat y}$ differ by at most $s_t$.
\end{lemma}
\begin{proof}
$\mathrm{gain}_t(x)=h_t+\sum_wN_{tw}x_w+\sum_{w\ne t}P_{tw}x_w$ and
$\sum_{w\ne t}P_{tw}x_w=\sum_k\lambda_kv_k(t)\,y_k(x^{t\to0})$. The gain under
$F'_{\hat y}$ replaces the last sum by $\sum_k\lambda_kv_k(t)\hat y_k$. Finally
$|y_k(x^{t\to0})-\hat y_k|\le|y_k(x)-\hat y_k|+|v_k(t)|$.
\end{proof}

\begin{lemma}[Exactness]\label{lem:exactC}
Every optimum $x^*$ is enumerated in a run whose grid point is within $\delta_k$
of $y_k(x^*)$ for all $k$.
\end{lemma}
\begin{proof}
Every flip at $x^*$ changes $F$ by at most zero, so by Lemma~\ref{lem:flip}
$\gamma_t(x^*)\le s_t$ for every creator on its path. No tie convention is
needed.
\end{proof}

\begin{theorem}[Round-based enumeration]\label{thm:rounds}
Let $D_N=\frac1n\sum_{u\ne v}N_{uv}^2>0$ and
$W_t=\sum_{s:\rho(s)<\rho(t)}N_{ts}^2$, and suppose
\begin{align}
 \|N\|^2&\le K\,D_N,\tag{C1}\label{C1}\\
 \sum_{t:\,2\le\rho(t)\le\ell}\sqrt{W_t}&\ge4\sqrt3\,g\,n\sqrt{D_N},\tag{C2}\label{C2}\\
 \sum_{t\le T}s_t&\le\tfrac18g\,n\sqrt{D_N},\tag{C3}\label{C3}
\end{align}
with $\|N\|$ the operator norm, and $n\ge256\pi\tau(\ell-1)K/g^2$. Put
$K'=(\ell-1)K$ and $\alpha=\max\{2,9\tau\sqrt{K'}/(4g)\}$. Then every run of
Algorithm~C enumerates at most $(6+2\ell)\,2^ne^{-c_*n}$ assignments, where
\[
 c_*=\min\Bigl\{\frac{g^2}{32\tau\alpha^2K'},\ \frac{g^2}{512\ell^2\tau K},\
 c_{\rm HW}\frac{g^2}{256\ell^2\tau K},\ c_{\rm HW}\frac{g}{16\ell\sqrt K}\Bigr\}.
\]
Algorithm~C is exact and runs in polynomial space and time
$\Ostar\bigl(|\mathcal Y|\,(2^{\tau n}+2^ne^{-c_*n})\bigr)$.
\end{theorem}

Condition~\eqref{C1} involves only the ordinary operator norm of the remainder,
not the norm of its triangular part as in the count-guessing note~\cite{CG}. The
price is that preferred values adapt to earlier rounds only.

\begin{lemma}[Race identity with rounds]\label{lem:race}
Let $x$ be uniform and $u=2x-1$. For a creator $t\in B_j$,
$\gamma_t(x)=|\Delta_t|-\frac12u_t\sum_{w\ne t:\rho(w)\ge j}N_{tw}u_w$. Hence $x$
passes~\eqref{eq:roundsum} if and only if $S'+Q_{\rm same}+Q_{\rm later}\le0$,
where
\begin{gather*}
 S'=\sum_{t\in I}(|\Delta_t|-s_t),\qquad
 Q_{\rm same}=-\frac12\sum_{j\le\ell}\ \sum_{t\in I\cap B_j}\
            \sum_{w\in B_j\setminus\{t\}}u_tN_{tw}u_w,\\
 Q_{\rm later}=-\frac12\sum_{w:\rho(w)\ge2}u_w\kappa_w,\qquad
 \kappa_w=\sum_{t\in I:\,\rho(t)<\rho(w)}u_tN_{tw}.
\end{gather*}
\end{lemma}
\begin{proof}
For $\rho(w)\ge j$ write $x_w=\frac12+\frac12u_w$; then the gain of $x_t$ under
$F'_{\hat y}$ is $\Delta_t+\frac12\sum_{w\ne t:\rho(w)\ge j}N_{tw}u_w$. For a
creator, $2p_t-1=-u_t$ and $(2p_t-1)\Delta_t=|\Delta_t|$.
\end{proof}

\begin{lemma}[Gap profile with rounds]\label{lem:profile}
Let $L$ be the $T\times T$ matrix with $L_{ts}=N_{ts}$ if $\rho(s)<\rho(t)$ and
$0$ otherwise, and $H_t=h^{\hat y}_t+\frac12\sum_{w\ne t}N_{tw}$. Then
$\Delta_t=H_t+\frac12(Lu)_t$, $\|L\|^2\le(\ell-1)\|N\|^2\le K'D_N$, and, with
$M=\alpha\sqrt{K'D_N}$,
\[
 \Prb\Bigl[\sum_{t\le T}\min\{|\Delta_t|,M\}<\tfrac34g\,n\sqrt{D_N}\Bigr]
 \le4e^{-g^2n/(16\tau K')}.
\]
\end{lemma}
\begin{proof}
The rows of $L$ in round $j$ form the block $N_{B_j\times B_{<j}}$ of norm at most
$\|N\|$, so $\|Lu\|^2=\sum_j\|N_{B_j\times B_{<j}}u\|^2\le(\ell-1)\|N\|^2\|u\|^2$.
The rest is the proof of the gap-profile lemma in~\cite{CG}, with $K'$ in place of
$K$. Split $t\le T$ into $\mathrm{Small}=\{|H_t|\le M/2\}$ and $\mathrm{Large}$.
For $t\in\mathrm{Small}$, $(a-b)_+\le a^2/(4b)$ gives
$\min\{|\Delta_t|,M\}\ge|\Delta_t|-(Lu)_t^2/(8M)$, a total loss of at most
$K'D_NT/(8M)$. For $t\in\mathrm{Large}$ with $|(Lu)_t|\le M/2$ the clipped gap is
at least $M/4$, and at most $4T/\alpha^2$ indices have $|(Lu)_t|>M/2$. Each
$(Lu)_t$ is a symmetric Rademacher sum with variance $W_t$, so by (T5)
$\E|\Delta_t|\ge\frac12\sqrt{W_t/3}$; also $M/4\ge\frac12\sqrt{W_t/3}$. By
\eqref{C2}, $\E G+\frac M4|\mathrm{Large}|\ge2gn\sqrt{D_N}$ for
$G=\sum_{\mathrm{Small}}|\Delta_t|$, while the two losses total at most
$\frac98T\sqrt{K'D_N}/\alpha\le\frac12gn\sqrt{D_N}$. The function $G$ is convex
with Lipschitz constant $\ell_G\le\frac12\sqrt{TK'D_N}$; (T3) with
$r=\frac12gn\sqrt{D_N}$, and $8\sqrt\pi\,\ell_G\le\frac14gn\sqrt{D_N}$ by the
bound on $n$, gives the claim.
\end{proof}

\begin{proof}[Proof of Theorem~\ref{thm:rounds}]
Exactness is Lemma~\ref{lem:exactC}. Fix a run and a uniform $x$.

\emph{Shortfall.} Let $Y_t=[t\in I]\min\{|\Delta_t|,M\}$. Given the variables of
the rounds before $j$, the values $\Delta_t$, $t\in B_j$, are fixed and the $u_t$,
$t\in B_j$, are independent and uniform, so the $Y_t$ are independent with values
in $[0,M]$ and means $\frac12\min\{|\Delta_t|,M\}$. By (T1),
$\sum_tY_t\ge\frac12\sum_t\min\{|\Delta_t|,M\}-a$ except with probability
$e^{-2a^2/(TM^2)}$. Take $a=\frac18gn\sqrt{D_N}$. With Lemma~\ref{lem:profile}
and~\eqref{C3}, $S'\ge\sum_tY_t-\sum_{t\le T}s_t\ge\frac18gn\sqrt{D_N}$ outside
two events.

\emph{Later interactions.} In round $j\ge2$, the increment of $Q_{\rm later}$ is
$-\frac12\sum_{w\in B_j}u_w\kappa_w$, where the $\kappa_w$ depend only on earlier
rounds and the $u_w$ are fresh. With $v_t=u_t[t\in I]$,
$\sum_{w\in B_j}\kappa_w^2=\|N_{B_{<j}\times B_j}^{\mathsf T}v\|^2\le\|N\|^2T$, so
$\sum_w\kappa_w^2\le\ell KD_NT$ on every path. By (T2) with $V=\frac14\ell KD_NT$,
$\Prb[Q_{\rm later}\le-b]\le e^{-2b^2/(\ell KD_NT)}$. Take
$b=\frac1{32}gn\sqrt{D_N}$.

\emph{Same-round interactions.} Fix $j$ and condition on the earlier rounds. Put
$c_t=2p_t-1$ and $\beta_t=-c_tu_t$ for $t\in B_j$; the $\beta_t$ are independent
uniform signs, $t\in I$ if and only if $\beta_t=1$, and $u_t=-c_t\beta_t$. With
the matrix $M^{(j)}_{tw}=c_tc_wN_{tw}$ on $B_j$, which has zero diagonal and
$\|M^{(j)}\|\le\|N\|$,
\[
 Q_{{\rm same},j}=-\frac12\sum_{t\ne w}\frac{1+\beta_t}2\,\beta_t\beta_wM^{(j)}_{tw}
 =-\frac14\beta^{\mathsf T}M^{(j)}\beta-\frac14\sum_w r_w\beta_w,\qquad
 r=M^{(j)}\one.
\]
Here $\|r\|^2\le\|N\|^2|B_j|$ and $\|M^{(j)}\|_F^2\le|B_j|\|N\|^2$. Let
$e=gn\sqrt{D_N}/(64\ell)$. By (T1) the linear part is below $-e$ with probability
at most $e^{-8e^2/(|B_j|KD_N)}$, and by (T4) the quadratic part is below $-e$ with
probability at most $\exp(-c_{\rm HW}\min\{16e^2/(|B_j|KD_N),\,4e/\sqrt{KD_N}\})$.
Summing over the rounds, $Q_{\rm same}\ge-\frac1{32}gn\sqrt{D_N}$ outside $2\ell$
events.

\emph{Conclusion.} Outside all exceptional events,
$S'+Q_{\rm same}+Q_{\rm later}\ge(\frac18-\frac1{32}-\frac1{32})gn\sqrt{D_N}>0$.
Using $|B_j|\le T\le\tau n$, the exponents of the exceptional probabilities are
$g^2n/(16\tau K')$, $g^2n/(32\tau\alpha^2K')$, $g^2n/(512\ell\tau K)$,
$g^2n/(512\ell^2\tau K)$, $c_{\rm HW}g^2n/(256\ell^2\tau K)$ and
$c_{\rm HW}gn/(16\ell\sqrt K)$, each at least $c_*n$ (the first because
$\alpha\ge2$). With the factor $4$ of Lemma~\ref{lem:profile}, their total
probability is at most $(6+2\ell)e^{-c_*n}$.
\end{proof}

\section{Max-Cut with arbitrary unit clauses on spectral expanders}
\label{sec:maxcut}

For a $d$-regular graph $G$ let $\lambda(G)=\max_{i\ge2}|\lambda_i(A_G)|$.

\begin{corollary}[Max-Cut on spectral expanders]\label{cor:maxcut}
Fix $\lambda_0\ge1$ and $0<\varepsilon\le\frac12$. There are $n_0$ and
$c=c(\lambda_0,\varepsilon)<2$ such that the following holds for every
$d$-regular graph $G$ on $n\ge n_0$ vertices with
\[
 64\lambda_0^2\le d\le(1-\varepsilon)n,\qquad \lambda(G)\le\lambda_0\sqrt d,
\]
and every multiset of unit clauses. The formula consisting of $(u\lor v)$ and
$(\neg u\lor\neg v)$ for every edge $uv$ of $G$, together with the unit clauses,
is solved by a deterministic algorithm in polynomial space and time
$\Ostar(c^n)$.
\end{corollary}

Without unit clauses this is Max-Cut on $G$, since the edge contributes
$1+[x_u\ne x_v]$. With unit clauses it is Max-Cut with arbitrary external fields.

\begin{proof}
Apply Theorem~\ref{thm:rounds} with $r=1$, $v_1=\one$, $\lambda_1=\beta:=-2d/(n-1)$,
exact counts $\mathcal Y_1=\{0,\dots,n\}$ and $\delta_1=0$, so $s_t=|\beta|$; with
$\tau=\frac12$, $\ell=2$, and $B_1,B_2$ the first and second $n/4$ variables of
any order. Here $J=-2A_G$ and $N=-2A_G-\beta(\one\one^{\mathsf T}-I)$.

\eqref{C1}: on $\one$, $N$ acts as $-2d-\beta(n-1)=0$; on $\one^\perp$ its
eigenvalues are $-2\lambda_i+\beta$. So $\|N\|\le2\lambda_0\sqrt d+2\le
3\lambda_0\sqrt d$. An edge has $N_{uv}=-2(1-d/(n-1))$ and a non-edge
$N_{uv}=-\beta$, so $D_N=4d(1-d/(n-1))\ge2\varepsilon d$ for $n\ge2/\varepsilon$.
Thus \eqref{C1} holds with $K=9\lambda_0^2/(2\varepsilon)$.

\eqref{C2}: since $1-d/(n-1)\ge\varepsilon/2$, an edge from $t\in B_2$ to $B_1$
contributes at least $\varepsilon^2$ to $W_t$. With $\deg_{B_1}(t)\le d$,
\[
 \sum_{t\in B_2}\sqrt{W_t}\ge\varepsilon\sum_{t\in B_2}\sqrt{\deg_{B_1}(t)}
 \ge\frac{\varepsilon\,e(B_1,B_2)}{\sqrt d}
 \ge\frac\varepsilon{\sqrt d}\Bigl(\frac{dn}{16}-\frac{\lambda(G)\,n}4\Bigr)
 \ge\frac{\varepsilon n\sqrt d}{32},
\]
by the expander mixing lemma
$|e(S,S')-d|S||S'|/n|\le\lambda(G)\sqrt{|S||S'|}$ for disjoint $S,S'$, and
$\sqrt d\ge8\lambda_0$. Since $\sqrt{D_N}\le2\sqrt d$, \eqref{C2} holds with
$g=\varepsilon/(256\sqrt3)$.

\eqref{C3}: $\sum_{t\le T}s_t=T|\beta|\le2d\le\frac18gn\sqrt{2\varepsilon d}$
once $\sqrt n\ge16\sqrt2/(g\sqrt\varepsilon)$, because $d\le n$.

All parameters depend only on $\lambda_0$ and $\varepsilon$, and $|\mathcal Y|=n+1$.
\end{proof}

\paragraph{Examples.}
\begin{itemize}
\item \emph{Paley graphs} $P_q$, $q\equiv1\pmod4$ prime: $d=(q-1)/2$ and the
nontrivial eigenvalues are $(-1\pm\sqrt q)/2$, so $\lambda_0=1$ and
$\varepsilon=\frac12$ work for $q\ge13$.
\item \emph{Random $d$-regular graphs}: $\lambda(G)=O(\sqrt d)$ with high
probability, for fixed $d$ by Friedman's theorem and for every
$d\le n/2$ by the result of Tikhomirov and Youssef~\cite{TY19}.
\item \emph{Ramanujan graphs} that are not bipartite, such as the
Lubotzky--Phillips--Sarnak graphs, with $\lambda(G)\le2\sqrt{d-1}$.
\end{itemize}
The numerical check in \texttt{explore\_paley\_maxcut.py} found the constant of
the triangular condition of~\cite{CG} to be about $1.3$ on Paley graphs in a
random order; Theorem~\ref{thm:rounds} makes that condition unnecessary.

\paragraph{Why no simpler algorithm is known for these classes.}
\begin{itemize}
\item Fixing all variables outside an independent set $S$ costs
$\Ostar(2^{n-|S|})$. By the Hoffman bound,
$\alpha(G)\le n\lambda(G)/(d+\lambda(G))\le\lambda_0n/\sqrt d$, so the saving
vanishes as $d$ grows.
\item Degree-based polynomial-space algorithms give
$\Ostar(2^{(1-2/\Delta)n})$ for Max-Cut~\cite{DKP07} and
$\Ostar(2^{n(1-\alpha\ln d/d)})$ for Max 2-CSP~\cite{GK}; both savings vanish as
$d$ grows.
\item Expanders have treewidth linear in $n$, so dynamic programming over
decompositions gives no saving, and also needs exponential space.
\item The automorphism groups of Paley and Cayley graphs have polynomial size and
give only polynomial savings.
\item Williams' algorithm~\cite{Williams} beats $2^n$ on every instance, but uses
exponential space. Whether general Max-Cut admits a polynomial-space algorithm
faster than $2^n$ is open~\cite{Lalovic24}.
\end{itemize}
We found no polynomial-space algorithm with a fixed saving for Max-Cut on these
families when $d$ grows; this search was not exhaustive. For $d<64\lambda_0^2$
the interaction graph has bounded degree, and polynomial-space algorithms for
Max 2-CSP of bounded degree~\cite{GK} give a fixed saving there, so the saving
can be made uniform over all $d\le(1-\varepsilon)n$.

\section{Count guessing without a given decomposition}\label{sec:spectral}

\begin{corollary}[Spectral count guessing]\label{cor:spectral}
Let $v_1,\dots,v_r$ be orthonormal and $\lambda_1,\dots,\lambda_r$ real, for
example approximate eigenpairs of $J$ of largest absolute value, and use the grids
$\mathcal Y_k=\delta\mathbb Z\cap[-\|v_k\|_1-\delta,\|v_k\|_1+\delta]$. Then
$|\mathcal Y|\le(2\sqrt n/\delta+3)^r$ and
$\sum_{t\le T}s_t\le\sum_k|\lambda_k|(\delta\sqrt n+1)$. Hence, if \eqref{C1}
and \eqref{C2} hold for $N$ and
$\sum_k|\lambda_k|(\delta\sqrt n+1)\le\frac18gn\sqrt{D_N}$, Max 2-SAT is solved in
polynomial space and time $\Ostar((2\sqrt n/\delta+3)^r(2^{\tau n}+2^ne^{-c_*n}))$.
\end{corollary}
\begin{proof}
$\sum_ts_t=\sum_k|\lambda_k|(\delta\|v_k\|_1+\|v_k\|_2^2)$ and
$\|v_k\|_1\le\sqrt n$. Exactness holds for any vectors (Lemma~\ref{lem:exactC}).
\end{proof}

The algorithm needs no planted structure as input: it computes approximate
eigenpairs of $J$ and checks the conditions. The next proposition shows that
random formulas whose parity means have low rank satisfy them.

\begin{proposition}[Planted low-rank parity structure]\label{prop:planted}
Let the coefficients $J_{uv}$, $u<v$, be independent with
$|J_{uv}-\E J_{uv}|\le B$ and $\nu\le\operatorname{Var}J_{uv}\le\chi\nu$, and
suppose $\E J_{uv}=\bar A_{uv}$ for $u\ne v$, where $\bar A$ is a symmetric matrix
of rank at most $r$ with $|\bar A_{uu}|\le B_0$ and $\|\bar A\|\le\mu n$. Let $P$
collect the eigenpairs of $J$ with $|\lambda|>\theta:=4\sqrt{\chi\nu n}$, and
split the prefix into two rounds of size $T/2$. There is $C=C(\chi,r,\tau,B_0)$
such that, if $n\nu\ge C(B^2\log n+\mu^2)$, then with probability $1-o(1)$
the conditions of Corollary~\ref{cor:spectral} hold with $K=O(\chi)$,
$g=\Omega(\tau^{3/2}/\sqrt\chi)$ and the grid spacing
$\delta=gn\sqrt{D_N}/(16\sqrt n\,\Lambda)$, $\Lambda=\sum_k|\lambda_k|$, which the
algorithm computes from the data. The running time is
$\Ostar(n^{O(r)}\,2^{(1-\gamma)n})$ with $\gamma>0$ depending only on $\chi$ and
$\tau$.
\end{proposition}
\begin{proof}
Let $X$ be the off-diagonal matrix of the $J_{uv}-\E J_{uv}$. By the norm bound of
Bandeira and van Handel~\cite{BvH16}, applied as in Proposition~12 of~\cite{CG},
$\|X\|\le3\sqrt{\chi\nu n}$ with probability $1-o(1)$. Then
$J=\bar A-\operatorname{diag}(\bar A)+X$ has at most $r$ eigenvalues of absolute
value above $\theta$ (Weyl), so $P$ has rank $r'\le r$, and $J-P$ has all
eigenvalues in $[-\theta,\theta]$. Hence $\bar A-P=(J-P)-X+\operatorname{diag}(\bar A)$
has rank at most $2r$ and norm at most $\theta+\|X\|+B_0$, and
\[
 N=X+\operatorname{offdiag}(\bar A-P),\qquad
 \|N\|\le\|X\|+2(\theta+\|X\|+B_0)=O(\sqrt{\chi\nu n}).
\]
The low-rank part has Frobenius norm at most
$\sqrt{2r}(\theta+\|X\|+B_0)=O(\sqrt{r\chi\nu n})$, against
$\|X\|_F\ge\sqrt{(1-\eta)n(n-1)\nu}$ by Bernstein's inequality, so
$D_N=(1-o(1))\frac1n\|X\|_F^2\ge(1-2\eta)(n-1)\nu$ and \eqref{C1} holds with
$K=O(\chi)$. For~\eqref{C2}, Bernstein's inequality gives
$\sum_{s\in B_1}X_{ts}^2\ge(1-\eta)|B_1|\nu$ for all $t\in B_2$, and the low-rank
part changes $\sum_{t\in B_2}\sqrt{W_t}$ by at most
$\sqrt{T}\cdot O(\sqrt{r\chi\nu n})=O(n\sqrt{r\chi\nu})$, which is negligible
against $|B_2|\sqrt{(1-\eta)|B_1|\nu}=\Theta(\tau^{3/2}n^{3/2}\sqrt\nu)$.
For~\eqref{C3}, the choice of $\delta$ gives $\Lambda\delta\sqrt n=
\frac1{16}gn\sqrt{D_N}$, so it remains to check $\Lambda\le\frac1{16}gn\sqrt{D_N}$.
Each of the at most $r$ selected eigenvalues has
$|\lambda_k|\le\|J\|\le\mu n+B_0+\|X\|$, and $\sqrt{D_N}\ge(1-o(1))\sqrt{n\nu}$;
so this holds once $n\nu\ge(32r/g)^2\mu^2$ and $n$ is large. The number of runs
is $(2\sqrt n/\delta+3)^r=(32\Lambda/(g\sqrt{D_N})+3)^r=n^{O(r)}$, since
$\Lambda/\sqrt{D_N}=O(r\mu\sqrt{n/\nu}+r\sqrt\chi)$ and $\nu\ge B^2\log n/n$.
\end{proof}

This removes the assumption ``given the decomposition'' from the count-guessing
theorem for planted bisections, planted group structure, and any other parity
pattern whose mean has bounded rank. Approximate eigenpairs suffice because
Lemma~\ref{lem:exactC} holds for any vectors and the conditions are checked for
the vectors actually used.

\section{What remains}\label{sec:remains}

\begin{itemize}
\item \emph{Complementary pairs} ($\pm J$ spin glasses on a fixed graph) need the
random-order version of Theorem~\ref{thm:weighted}.
\item \emph{Irregular expanders.} Theorem~\ref{thm:rounds} measures $\|N\|$ against
the average $D_N$; for very irregular degree sequences a weighted version, as in
Theorem~\ref{thm:weighted}, is the natural next step.
\item \emph{Biased random formulas at density} $\omega(1)\le d<C\log n$, by
branching first on the few vertices of large degree.
\item \emph{Aligned structure of high rank}, such as noisy blow-ups, where the
counts of many small groups would have to be guessed.
\item The literature search behind Section~\ref{sec:maxcut} should be extended
before any claim of novelty.
\end{itemize}
The script \texttt{explore\_weighted\_rounds.py} checks exactness of Algorithms
A and~C, including the spectral grid, against exhaustive search: no optimum was
lost on 200 formulas for Algorithm~A (which enumerated $8.6\%$ of all assignments
on average), nor on 120 formulas each for exact counts and spectral grids in
Algorithm~C ($24\%$ and $21\%$ enumerated per run). These small instances test
correctness only, not the asymptotic bounds.

\begin{thebibliography}{99}
\small\raggedright
\bibitem{BvH16} A.~S.~Bandeira and R.~van Handel. Sharp nonasymptotic bounds on
 the norm of random matrices with independent entries. \emph{Ann. Probab.}
 44(4):2479--2506, 2016.
\bibitem{CG} \emph{Count guessing makes constraints useful}, working note,
 11 October 2026 (\texttt{max2sat\_count\_guessing.tex}).
\bibitem{DKP07} F.~Della~Croce, M.~J.~Kaminski and V.~Th.~Paschos. An exact
 algorithm for MAX-CUT in sparse graphs. \emph{Oper. Res. Lett.} 35(3):403--408,
 2007.
\bibitem{EE} \emph{Max 2-SAT by early constraint-guided enumeration}, working
 note, 8 October 2026 (\texttt{max2sat\_early\_enumeration.tex}).
\bibitem{GK} A.~Golovnev and K.~Kutzkov. New exact algorithms for the
 2-constraint satisfaction problem. \emph{Theoret. Comput. Sci.} 526:18--27, 2014.
\bibitem{Lalovic24} K.~Lalovic. Exact algorithms for MaxCut on split graphs.
 arXiv:2405.20599, 2024.
\bibitem{Paper} \emph{Dense Max 2-SAT in polynomial space}, research draft,
 12 October 2026 (\texttt{dense\_max2sat\_paper.tex}).
\bibitem{RV13} M.~Rudelson and R.~Vershynin. Hanson--Wright inequality and
 sub-gaussian concentration. \emph{Electron. Commun. Probab.} 18(82):1--9, 2013.
\bibitem{Tal95} M.~Talagrand. Concentration of measure and isoperimetric
 inequalities in product spaces. \emph{Publ. Math. IH\'ES} 81:73--205, 1995.
\bibitem{TY19} K.~Tikhomirov and P.~Youssef. The spectral gap of dense random
 regular graphs. \emph{Ann. Probab.} 47(1):362--419, 2019.
\bibitem{Williams} R.~Williams. A new algorithm for optimal 2-constraint
 satisfaction and its implications. \emph{Theoret. Comput. Sci.}
 348(2--3):357--365, 2005.
\end{thebibliography}
\end{document}
